%%%%%%%%%%%%%%%%%%%%%%%%%%%%%%%%%%%%%%%%%%%%%%%%%%%%%%%%%%%%%%%%%%%%%%%%%%%%%%%%
\chapter{Relativistic Diffusion Processes:
Existing Approaches and Their Relation to SHP}
\label{ch:RelativisticDiffusions}
%%%%%%%%%%%%%%%%%%%%%%%%%%%%%%%%%%%%%%%%%%%%%%%%%%%%%%%%%%%%%%%%%%%%%%%%%%%%%%%%
\section{Introduction}
The construction of a relativistic analogue of Brownian motion has been
an active research topic for more than sixty years. Although Brownian
motion is one of the fundamental stochastic processes in classical
probability theory, its extension to relativistic physics is far from
straightforward.
The difficulty does not arise from stochastic calculus itself but from
the geometric structure of spacetime. Ordinary Brownian motion is
defined on Euclidean space with an absolute time parameter. In
relativistic mechanics there is no preferred time coordinate, velocities
are constrained by the speed of light, and the state space is naturally
described by Lorentzian geometry rather than Euclidean geometry.
Consequently, many apparently natural generalizations of Brownian motion
either violate Lorentz covariance, allow superluminal propagation, or
fail to possess satisfactory probabilistic properties.
During the past several decades numerous relativistic diffusion models
have been proposed. These approaches differ substantially in their
mathematical formulation, physical motivation, and intended domain of
application. Some are formulated directly on Minkowski spacetime, others
on the mass-shell hyperboloid, while still others employ the tangent
bundle or orthonormal frame bundle of spacetime.
The objective of this chapter is not to advocate a particular model but
to provide a systematic comparison of the principal approaches that have
appeared in the literature. Special attention will be devoted to those
features that are relevant for a stochastic extension of the
Stueckelberg--Horwitz--Piron (SHP) framework developed in subsequent
chapters.
Throughout this chapter we distinguish carefully between three
conceptually different questions:
\begin{enumerate}
\item
How can one construct a Lorentz-covariant stochastic process?
\item
What is the physically appropriate state space for relativistic
diffusion?
\item
How should the stochastic evolution be parametrized?
\end{enumerate}
As we shall see, different authors provide different answers to each of
these questions.
One conclusion will become apparent from the outset. There is no unique
relativistic analogue of Brownian motion. Instead, there exists a family
of mathematically consistent diffusion processes, each adapted to a
particular physical interpretation. This observation suggests that the
appropriate stochastic process for SHP should be derived from the
structure of SHP itself rather than imported directly from
nonrelativistic stochastic mechanics.
%%%%%%%%%%%%%%%%%%%%%%%%%%%%%%%%%%%%%%%%%%%%%%%%%%%%%%%%%%%%%%%%%%%%%%%%%%%%%%%%
\section{Why Relativistic Diffusion is Difficult}
\label{sec:WhyRelativisticDiffusion}
%%%%%%%%%%%%%%%%%%%%%%%%%%%%%%%%%%%%%%%%%%%%%%%%%%%%%%%%%%%%%%%%%%%%%%%%%%%%%%%%
At first sight, constructing a relativistic version of Brownian motion
appears to be a straightforward task. One might attempt simply to
replace the spatial coordinates of nonrelativistic diffusion by
space-time coordinates and require Lorentz covariance of the resulting
equations.
This seemingly natural procedure fails for several fundamental reasons.
The difficulties are not merely technical. They arise from the geometric
structure of Minkowski space-time and from the conceptual foundations of
both probability theory and relativity. In this section we examine these
obstacles individually, since they motivate nearly every relativistic
diffusion model discussed later in this chapter.
%%%%%%%%%%%%%%%%%%%%%%%%%%%%%%%%%%%%%%%%%%%%%%%%%%%%%%%%%%%%%%%%%%%%%%%%%%%%%%%%
\subsection{Brownian Motion Requires an Ordering Parameter}
%%%%%%%%%%%%%%%%%%%%%%%%%%%%%%%%%%%%%%%%%%%%%%%%%%%%%%%%%%%%%%%%%%%%%%%%%%%%%%%%
The standard Wiener process
\[W(t)\]
is defined with respect to a monotonically increasing parameter
\(t\), satisfying
\[W(0)=0,\]
together with the increment property
\[W(t+\Delta t)-W(t)\sim N(0,\Delta t).\]
The existence of such an ordering parameter is essential.
Successive increments are defined only after the stochastic trajectory
has been ordered.
In Newtonian mechanics this causes no difficulty because absolute time
provides a universal ordering variable.
Special relativity, however, contains no preferred global time.
Different inertial observers assign different coordinate times to the
same physical events.
Consequently, coordinate time cannot serve as a fundamental stochastic
parameter without sacrificing manifest covariance.
This is perhaps the first conceptual obstacle encountered by every
relativistic diffusion theory.
%%%%%%%%%%%%%%%%%%%%%%%%%%%%%%%%%%%%%%%%%%%%%%%%%%%%%%%%%%%%%%%%%%%%%%%%%%%%%%%%
\subsection{Coordinate Time Is Observer Dependent}
%%%%%%%%%%%%%%%%%%%%%%%%%%%%%%%%%%%%%%%%%%%%%%%%%%%%%%%%%%%%%%%%%%%%%%%%%%%%%%%%
Suppose one writes
\begin{equation}
dx^\mu=b^\mu dt+\sigma^\mu{}_a\,dW^a(t),
\label{eq:CoordinateTimeDiffusion}
\end{equation}
where \(t=x^0/c\).
Although mathematically well defined in a chosen inertial frame,
Equation~(\ref{eq:CoordinateTimeDiffusion}) is not manifestly Lorentz
covariant.
After a Lorentz transformation,
\[t\rightarrow t',\]
the Wiener process itself changes parametrization.
The transformed stochastic process is generally no longer a Wiener
process with respect to \(t'\).
Therefore,
ordinary Brownian motion is not preserved by Lorentz transformations.
This simple observation explains why relativistic stochastic mechanics
cannot be obtained merely by replacing spatial vectors with four-vectors.
%%%%%%%%%%%%%%%%%%%%%%%%%%%%%%%%%%%%%%%%%%%%%%%%%%%%%%%%%%%%%%%%%%%%%%%%%%%%%%%%
\subsection{The Speed-of-Light Constraint}
%%%%%%%%%%%%%%%%%%%%%%%%%%%%%%%%%%%%%%%%%%%%%%%%%%%%%%%%%%%%%%%%%%%%%%%%%%%%%%%%
A second difficulty arises from the kinematics of Brownian motion.
In ordinary diffusion,
increments satisfy
\[\Delta x\sim\sqrt{\Delta t}.\]
The corresponding instantaneous velocity behaves formally as
\[\frac{\Delta x}{\Delta t}\sim\frac{1}{\sqrt{\Delta t}},\]
which diverges as
\[\Delta t\rightarrow0.\]
Thus Brownian trajectories are almost surely nowhere differentiable.
In nonrelativistic mechanics this presents no contradiction because
there is no upper limit on velocity.
In relativity,
however,
physical trajectories must satisfy
\[|\mathbf v|<c.\]
Consequently,
ordinary Wiener trajectories cannot represent physical particle
worldlines.
One must either smooth the trajectories or formulate diffusion in a
different state space.
%%%%%%%%%%%%%%%%%%%%%%%%%%%%%%%%%%%%%%%%%%%%%%%%%%%%%%%%%%%%%%%%%%%%%%%%%%%%%%%%
\subsection{The Geometry of Velocity Space}
%%%%%%%%%%%%%%%%%%%%%%%%%%%%%%%%%%%%%%%%%%%%%%%%%%%%%%%%%%%%%%%%%%%%%%%%%%%%%%%%
Relativistic velocities are not elements of Euclidean space.
For massive particles,
the four-velocity satisfies
\begin{equation}
u^\mu u_\mu=c^2.
\label{eq:MassShell}
\end{equation}
The allowed velocities therefore lie on the future hyperboloid
\[H^3=\{u^\mu:u^\mu u_\mu=c^2,\;u^0>0\}.\]
This hyperboloid possesses constant negative curvature.
Consequently,
relativistic diffusion naturally becomes diffusion on a curved manifold
rather than in ordinary Euclidean space.
This geometric observation forms the basis of Dudley's construction and
many later developments.
%%%%%%%%%%%%%%%%%%%%%%%%%%%%%%%%%%%%%%%%%%%%%%%%%%%%%%%%%%%%%%%%%%%%%%%%%%%%%%%%
\subsection{Which Variables Should Diffuse?}
%%%%%%%%%%%%%%%%%%%%%%%%%%%%%%%%%%%%%%%%%%%%%%%%%%%%%%%%%%%%%%%%%%%%%%%%%%%%%%%%
In nonrelativistic mechanics,
configuration-space diffusion is usually sufficient because velocity may
be obtained by differentiation.
Relativistically,
this distinction becomes more subtle.
Several possibilities exist.
\begin{enumerate}
\item
Diffusion in space-time.
\item
Diffusion of four-velocity.
\item
Diffusion in momentum space.
\item
Diffusion on the tangent bundle.
\item
Diffusion on phase space.
\item
Diffusion on the orthonormal frame bundle.
\end{enumerate}
Each choice leads to a different stochastic process.
Moreover,
the physically appropriate choice depends upon the intended
interpretation.
For example,
kinetic theory naturally favors momentum-space diffusion,
whereas stochastic mechanics emphasizes trajectories in configuration
space.
%%%%%%%%%%%%%%%%%%%%%%%%%%%%%%%%%%%%%%%%%%%%%%%%%%%%%%%%%%%%%%%%%%%%%%%%%%%%%%%%
\subsection{Proper Time Versus Coordinate Time}
%%%%%%%%%%%%%%%%%%%%%%%%%%%%%%%%%%%%%%%%%%%%%%%%%%%%%%%%%%%%%%%%%%%%%%%%%%%%%%%%
The invariant proper time
\[d\tau^2=dt^2-\frac{d\mathbf x^2}{c^2}\]
appears at first sight to solve the ordering problem.
Indeed,
many relativistic diffusion models employ proper time as the evolution
parameter.
However,
proper time itself introduces new questions.
For massless particles,
proper time is undefined.
For interacting many-particle systems,
individual particles possess different proper times.
Furthermore,
in quantum theory the notion of a sharply defined proper-time trajectory
is itself problematic.
The SHP framework avoids many of these difficulties by introducing an
independent invariant evolution parameter
\[\tau,\]
which exists even for off-shell motion.
This distinction will become increasingly important in later chapters.
%%%%%%%%%%%%%%%%%%%%%%%%%%%%%%%%%%%%%%%%%%%%%%%%%%%%%%%%%%%%%%%%%%%%%%%%%%%%%%%%
\subsection{Probability Density}
%%%%%%%%%%%%%%%%%%%%%%%%%%%%%%%%%%%%%%%%%%%%%%%%%%%%%%%%%%%%%%%%%%%%%%%%%%%%%%%%
In Euclidean diffusion,
the probability density satisfies
\[\rho(\mathbf x,t)\ge0,\]
together with
\[\int\rho(\mathbf x,t)\,d^3x=1.\]
Relativistically,
the corresponding probability interpretation is much less obvious.
Should one normalize over
\[d^3x,\]
over
\[d^4x,\]
over the mass shell,
or over phase space?
Different relativistic diffusion theories answer this question
differently.
The associated Fokker--Planck equations therefore differ as well.
%%%%%%%%%%%%%%%%%%%%%%%%%%%%%%%%%%%%%%%%%%%%%%%%%%%%%%%%%%%%%%%%%%%%%%%%%%%%%%%%
\subsection{Markovian Versus Non-Markovian Dynamics}
%%%%%%%%%%%%%%%%%%%%%%%%%%%%%%%%%%%%%%%%%%%%%%%%%%%%%%%%%%%%%%%%%%%%%%%%%%%%%%%%
Most relativistic diffusion models assume Markovian evolution,
primarily because it simplifies mathematical analysis.
However,
as discussed in the previous chapter,
both stochastic electrodynamics and generalized Langevin theory suggest
that realistic stochastic dynamics may possess finite correlation times
and memory effects.
From this perspective,
Markovian relativistic diffusion should be regarded as the simplest
approximation rather than the most general possibility.
Whether a covariant non-Markovian theory can be constructed remains an
important open problem.
%%%%%%%%%%%%%%%%%%%%%%%%%%%%%%%%%%%%%%%%%%%%%%%%%%%%%%%%%%%%%%%%%%%%%%%%%%%%%%%%
\subsection*{Critical Assessment}
%%%%%%%%%%%%%%%%%%%%%%%%%%%%%%%%%%%%%%%%%%%%%%%%%%%%%%%%%%%%%%%%%%%%%%%%%%%%%%%%
The principal obstacle in relativistic diffusion is not the definition
of random variables but the simultaneous reconciliation of stochastic
calculus with Lorentz geometry. The absence of a preferred time
parameter, the boundedness of physical velocities, the hyperbolic
geometry of the mass shell, and the ambiguity of the appropriate state
space collectively distinguish relativistic diffusion from its
nonrelativistic counterpart.
The diversity of approaches found in the literature reflects different
choices regarding these foundational issues rather than simple
differences in mathematical technique. Some formulations prioritize
manifest covariance, others emphasize kinetic theory or statistical
mechanics, while still others are designed for stochastic mechanics or
quantum foundations.
From the standpoint of the present monograph, perhaps the most important
observation is that the SHP framework already provides one of the key
ingredients absent from ordinary relativistic mechanics: an invariant
evolution parameter. This does not by itself solve the problem of
constructing a relativistic diffusion process, but it removes one of the
principal conceptual obstacles. As a result, SHP offers a particularly
natural setting in which to revisit Nelson's stochastic construction.
